% section_1_intro.tex — v6 from-scratch rewrite
% Spec: docs/Draft/THESIS_SKELETON.md v6 §1 + project_dwz_anchored_framing_locked_2026_04_27.md
% Numbers: docs/Draft/CANONICAL_FACT_SHEET.md
% Headline contribution sentence required in Contributions §1

\section{Introduction}
\label{sec:intro}

\subsection{Motivation}

Quarterly earnings calls offer a higher-frequency window into managerial communication than the annual 10-K disclosure cycle, and a richer one in source-identifiable speaker attribution: a single call admits separate prepared-remarks and spontaneous Q\&A segments and identifies who is speaking (CEO, CFO, other manager, analyst). \citeA{dzielinski2021} show that the prepared and spontaneous segments carry systematically different uncertainty content, and that within the Q\&A segment one can further isolate a persistent CEO-specific style component (\texttt{ClarityCEO}, recovered as the negative of a CEO fixed effect) from a call-level state residual (\texttt{UncResCEO}). Their three-component architecture is the starting point for this paper.

\citeauthor{dzielinski2021} validate their three-component decomposition on a nine-outcome set spanning market reactions (short- and long-window cumulative abnormal returns, abnormal trading volume), analyst behavior (forecast revisions, median recommendations), firm performance and valuation (Tobin's Q, return on assets), and governance (CEO compensation), as enumerated in their Appendix Table~A.1. Their outcome set does not include firm-level financing-policy decisions. The natural next question---and the central question of this paper---is whether the same three-component speech-uncertainty signal carries information for the firm's own \emph{financing} decisions: cash holdings, the precautionary liquidity buffer formalized by \citeA{opler1999} and updated for the post-2000 sample by \citeA{bates2009}.

\subsection{Gap}

To our knowledge, no prior work has applied a CEO-specific three-component speech-uncertainty decomposition to firm-quarter cash holdings, nor jointly examined the firm's own cash-financing margin and the outside-world's reaction margin (post-call quoted spread, regulatory comment-letter receipt) of CEO speech uncertainty under a common decomposition. \citeauthor{dzielinski2021}'s outcome set spans market reactions, analyst behavior, firm performance and valuation, and CEO compensation; the firm-side cash response is conceptually distinct (insider financing decisions versus outsider price formation and analyst behavior) and theoretically anchored elsewhere (precautionary motive of \citeauthor{opler1999}). The CEO-specific aggregation we adopt is theoretically motivated by \citeauthor{dzielinski2021}'s scripted-versus-improvised distinction between the prepared-remarks and Q\&A segments, together with the manager-fixed-effect-on-cash channel of \citeA{bertrand2003}: both anchor the cash-policy response to the CEO speaker specifically rather than to a manager-pool aggregate.

\subsection{Approach}

We preserve the \citeauthor{dzielinski2021} decomposition. The Q\&A speech-uncertainty share is recovered using the \citeA{loughran2011} finance-specific uncertainty word list at speaker-utterance resolution, aggregated over CEO speakers in the call's Q\&A segment. We then estimate a CEO fixed-effect regression of this share on Bates-style controls and recover the negative of the CEO fixed effect as \texttt{ClarityCEO} (a CEO-level trait following \citeA{bertrand2003}, who establish that manager fixed effects significantly predict corporate cash holdings) and the regression residual as \texttt{UncResCEO} (a call-level state residual). The presentation-segment analogue \texttt{UncPreCEO} is computed analogously on prepared remarks. We test the cash response (HC, \S\ref{sec:main:hc}) and its credit-constraint moderation (HFC, \S\ref{sec:main:hfc}) on a four-step firm/industry $\times$ year/year-quarter fixed-effects ladder of $112{,}968$ earnings-call transcripts covering $2{,}429$ U.S.\ firms over fiscal years 2002--2018, sourced from S\&P Capital~IQ and merged with Compustat quarterly fundamentals and CRSP daily security data. Inference on the speech-uncertainty regressors is one-tailed in the theoretically predicted direction \cite{opler1999,bertrand2003,almeida2004}.

\subsection{Empirical Setting and Findings}

\textbf{Cash response.} The Q\&A residual \texttt{UncResCEO} is positive and statistically significant at the ten-percent threshold in all twelve specifications of suite H1 (eight of twelve at the five-percent threshold in the Full method; twelve of twelve at the ten-percent threshold and seven of twelve at the five-percent threshold in the QtrExp expanding-window method), with point estimates between $+0.0016$ and $+0.0028$ in Full and $+0.0016$ and $+0.0046$ in QtrExp. The persistent CEO-style trait \texttt{ClarityCEO} is negative and significant in nine of twelve specifications under both methods, with point estimates between $-0.0170$ and $-0.0046$ in Full and $-0.0215$ and $-0.0024$ in QtrExp; the QtrExp method strengthens the persistent-style loading to nine of twelve at the five-percent threshold and six of twelve at the one-percent threshold. The presentation-segment share \texttt{UncPreCEO} is null across the same panel. The headline pattern of this paper is that the three-component speech-uncertainty decomposition \citeauthor{dzielinski2021} introduce carries information for firm-level cash policy: \emph{applied to a previously-untested outcome class, the same constructs deliver systematic and pre-registered loadings on the firm's precautionary cash buffer}.

\textbf{Constraint moderation.} Among the rated subsample with available Compustat credit-rating coverage (2002--2016), the \texttt{UncResCEO}~$\times$~Unrated interaction is positive at the five-percent threshold in two of eight cells on the one-quarter-lead cash dependent variable, with the credit-constrained Unrated firms of \citeA{faulkender2006} amplifying the residual response by approximately three-fold over the rated pool (HFC, \S\ref{sec:main:hfc}). A second constraint-channel moderator constructed on \citeauthor{hanqiu2007}'s sixteen-quarter cash-flow-volatility measure produces the same lead-dependent-variable amplification pattern at one-tier-weaker statistical significance (\S\ref{sec:main:cfvol}); the two constraint operationalizations \cite{faulkender2006,hanqiu2007} agree on the channel direction.

\textbf{Driver-side construct validation.} The CEO speech-uncertainty metric itself responds positively to four independently measured uncertainty drivers (\S\ref{sec:additional:drivers}): the firm-level political-risk index of \citeA{hassan2020}, the US news-based economic-policy-uncertainty index of \citeA{baker2016}, the global extension of \citeA{davis2016}, and product-market competition pressure measured by the text-based similarity index of \citeA{hoberg2016}.

\textbf{Outside-world reaction.} The post-call closing-quote relative bid-ask spread averaged over the 25-trading-day window of \citeA{bushee2018} (\S\ref{sec:additional:reaction:spread}) and the within-quarter SEC comment-letter receipt indicator of \citeA{lerman2026} (\S\ref{sec:additional:reaction:sec}) both load the \emph{presentation-segment} uncertainty component \texttt{UncPreCEO}---4 of 12 spread cells and 4 of 6 SEC cells significant at $p<0.10$ one-tailed---with the persistent style and Q\&A residual components null in both outside-world regressions. The asymmetry sits cleanly on \citeauthor{dzielinski2021}'s own theoretical structure: the segment outsiders react to is the one \citeauthor{dzielinski2021} characterize as ``scripted'' and ``firm-culture-reflecting,'' while the segments that drive the firm's cash decisions (Q\&A persistent and Q\&A residual) are the ones they characterize as ``revealing CEO style.''

\textbf{Endogeneity.} Section~\ref{sec:additional:endo} reports three identification designs against three orthogonal threats. A CEO sudden-death difference-in-differences (\S\ref{sec:additional:endo:death}) identifies a reverse-causality-resistant exogenous shock at small but cleanly-identified treated-firm count. A first-difference replication of \citeauthor{dzielinski2021}'s \S6 specification around CEO turnover (\S\ref{sec:additional:endo:dwzfd}) addresses time-invariant firm and manager heterogeneity. A \citeA{lewbel2012} heteroskedasticity-based instrumental-variable estimator (\S\ref{sec:additional:endo:lewbel}) addresses time-varying omitted confounders. The three designs cover orthogonal identification threats; we frame this as a three-threats package rather than as convergence on a single point estimate.

\subsection{Contributions}

The paper makes three contributions.

\begin{enumerate}
\item \textbf{Extending the DWZ three-component decomposition to corporate financing-policy outcomes.} \citeA{dzielinski2021} introduce and validate the persistent-trait, Q\&A-residual, presentation-share decomposition on market reactions, analyst forecasts and recommendations, firm value and performance, and CEO compensation; their paper does not test firm-level financing-policy outcomes. We apply their three-component architecture to the canonical precautionary-cash response of \citeA{opler1999} and \citeA{bates2009} and document that the Q\&A residual \texttt{UncResCEO} loads positively on cash in all twelve fixed-effects specifications (\S\ref{sec:main:hc}). The decomposition's three components carry pre-registered, theoretically grounded signal for firm-internal financing decisions in addition to the market and governance outcomes \citeauthor{dzielinski2021} document.

\item \textbf{Insider-outsider channel asymmetry on a common decomposition.} The same three-component decomposition that produces the cash-side result reproduces an insider-outsider channel asymmetry on the outsider-reaction outcomes (\S\ref{sec:additional:reaction}). Insiders' cash decisions track Q\&A components (\texttt{ClarityCEO} and \texttt{UncResCEO}); the post-call spread and SEC comment-letter outsiders track the presentation-segment component (\texttt{UncPreCEO}). The asymmetry is not imposed; it emerges within \citeauthor{dzielinski2021}'s own scripted-versus-improvised theoretical structure.

\item \textbf{Three-threats endogeneity package.} Three identification designs cover orthogonal threats to identification of the speech-uncertainty effect on cash (\S\ref{sec:additional:endo}). The package is presented as covering three threats with three designs, not as convergence on a single point estimate.
\end{enumerate}

\subsection{Roadmap}

Section~\ref{sec:framework} develops the conceptual framework and pre-commitment statement (\S\ref{sec:framework:precommit}), the speech-uncertainty measurement architecture (\S\ref{sec:framework:measurement}), the precautionary-motive and constraint-amplification literature anchors for HC and HFC (\S\ref{sec:framework:precautionary}), and the empirical design (\S\ref{sec:framework:design}). Section~\ref{sec:main} reports the data, sample, variables, and specification (\S\ref{sec:main:setup}); the cash-holdings result (HC, \S\ref{sec:main:hc}); the credit-constraint moderation (HFC, \S\ref{sec:main:hfc}); and the cash-flow-volatility moderator (\S\ref{sec:main:cfvol}). Section~\ref{sec:additional} reports the four driver regressions (\S\ref{sec:additional:drivers}), the post-call spread reaction (\S\ref{sec:additional:reaction:spread}), the SEC comment-letter reaction (\S\ref{sec:additional:reaction:sec}), and the three-threats endogeneity package (\S\ref{sec:additional:endo}). Section~\ref{sec:conclusion} synthesizes the channel asymmetries, discusses limitations, and outlines future research.
